% ======================================================================
%  "ONE" -- Auburn University research poster, third variant
%
%  One clear mathematical idea placed in a quiet, beautiful space.
%
%  Design language (the "Steve Jobs direction", 2026-07-01):
%    * warm paper-white ground (#F7F4EE), deep charcoal ink (#1E1E1E),
%      soft gray for secondary text (#5C5C5C), ONE quiet blue-gray
%      accent (#3F5E7A), a whisper of warm beige (#C7B8A1) in hairlines;
%    * two typefaces only: EB Garamond (title, theorem, takeaway) and
%      Inter (body), with STIX Two Math for mathematics;
%    * three asymmetric columns -- the centre is wider, more open, and
%      set slightly larger; the sides are quiet and explanatory;
%    * the main theorem sits in one soft, thin-bordered panel
%      (#E9EEF2 on #D8E1E8) -- no shadows, no loud boxes;
%    * one strong visual: a projected slice of the Fermat quintic
%      Calabi-Yau threefold, rendered in the poster's own palette.
%
%  120 x 72 cm landscape.  Compiles to a tagged, PDF/UA-1 PDF.
%
%  COMPILE WITH LuaLaTeX:
%    * Overleaf:  Menu -> Compiler -> LuaLaTeX, then Recompile.
%    * Locally:   latexmk poster-one.tex
%
%  Siblings: poster-research.tex (navy/orange, sheaf theory) and
%  poster-research-classic.tex (beamer, three equal columns).
% ======================================================================
\DocumentMetadata{lang=en-US, pdfstandard=ua-1, tagging=on}
\documentclass[12pt]{article}

% Fail fast, with a clear message, if the wrong engine is used.
\usepackage{iftex}
\RequireLuaTeX

% ----- page: 120 x 72 cm landscape, generous margins -------------------
\usepackage[paperwidth=120cm, paperheight=72cm,
            top=2.6cm, bottom=2.2cm, left=4.0cm, right=4.0cm]{geometry}
\pagestyle{empty}
\setlength{\parindent}{0pt}
\setlength{\parskip}{0.55em}
\setlength{\emergencystretch}{1.5em}  % kill loose justified lines

% ----- two typefaces: EB Garamond (display) + Inter (text) -------------
% Both are bundled in fonts/ (SIL OFL), so Overleaf renders identically.
% Inter ships no italic; emphasis in sans runs uses \auterm (semibold),
% and italics live only in the Garamond runs.
\usepackage{fontspec}
\setmainfont{Inter}[Path=fonts/, Extension=.otf,
  UprightFont=*-Regular, BoldFont=*-Bold]
\newfontfamily\SansMedium{Inter}[Path=fonts/, Extension=.otf,
  UprightFont=*-Medium, Ligatures=TeX]
\newfontfamily\SansSemibold{Inter}[Path=fonts/, Extension=.otf,
  UprightFont=*-SemiBold, Ligatures=TeX]
\newfontfamily\Serif{EBGaramond}[Path=fonts/, Extension=.otf,
  UprightFont=*-Regular, ItalicFont=*-Italic, BoldFont=*-Bold,
  Ligatures=TeX]
\newfontfamily\SerifMedium{EBGaramond}[Path=fonts/, Extension=.otf,
  UprightFont=*-Medium, Ligatures=TeX]
\newfontfamily\SerifSemibold{EBGaramond}[Path=fonts/, Extension=.otf,
  UprightFont=*-SemiBold, Ligatures=TeX]
\usepackage{mathtools}
\usepackage{booktabs}
\usepackage{graphicx}
\usepackage{tikz}
\usepackage{unicode-math}
\setmathfont{STIXTwoMath-Regular.otf}[Path=fonts/]
\setoperatorfont\symup

% ----- the colour system ------------------------------------------------
\usepackage{xcolor}
\definecolor{PaperBG}{HTML}{F7F4EE}     % warm paper-white ground
\definecolor{InkMain}{HTML}{1E1E1E}     % deep charcoal: body, formulas
\definecolor{InkSoft}{HTML}{5C5C5C}     % secondary: captions, refs
\definecolor{AccentBlue}{HTML}{3F5E7A}  % ONE accent: headings, labels
\definecolor{ThmBG}{HTML}{E9EEF2}       % theorem panel ground
\definecolor{ThmBorder}{HTML}{D8E1E8}   % theorem panel hairline
\definecolor{DividerWarm}{HTML}{D8D2C8} % light dividers
\definecolor{DiagSoft}{HTML}{8A99A8}    % diagram secondary
\definecolor{WarmAccent}{HTML}{C7B8A1}  % optional warm hairline
\colorlet{WarmDeep}{WarmAccent!55!InkMain} % readable warm, same family
\pagecolor{PaperBG}
\color{InkMain}

% ----- tagged figures (hero image, diagram, QR), thesis pattern --------
\usepackage{qrcode}
\usepackage{zref-savepos}
\ExplSyntaxOn
\box_new:N \l__pfig_box
\int_new:N \g__pfig_int
\tl_new:N  \l__pfig_alt_tl
\cs_generate_variant:Nn \tag_struct_begin:n { e }
% \PosterFigure{<alt text>}{<content>} -- tag <content> as ONE Figure
% with alt text and a bounding box, exactly like an included image.
\NewDocumentCommand{\PosterFigure}{mm}{%
  \leavevmode
  \tl_set:Ne \l__pfig_alt_tl {#1}%
  \tag_struct_begin:e { tag=Figure , alt={\tl_to_str:N \l__pfig_alt_tl} }%
  \tag_mc_begin:n {}%
  \tag_suspend:n {pfig}%
  \hbox_set:Nn \l__pfig_box {#2}%
  \tag_resume:n {pfig}%
  \int_gincr:N \g__pfig_int
  \zsavepos { pfig \int_use:N \g__pfig_int }%
  \tag_struct_gput:ene { \tag_get:n {struct_num} } { attribute }
    { /O~/Layout~/BBox~[
        \dim_to_decimal_in_bp:n { \zposx{pfig\int_use:N\g__pfig_int} sp } ~
        \dim_to_decimal_in_bp:n { \zposy{pfig\int_use:N\g__pfig_int} sp
                                  - \box_dp:N \l__pfig_box } ~
        \dim_to_decimal_in_bp:n { \zposx{pfig\int_use:N\g__pfig_int} sp
                                  + \box_wd:N \l__pfig_box } ~
        \dim_to_decimal_in_bp:n { \zposy{pfig\int_use:N\g__pfig_int} sp
                                  + \box_ht:N \l__pfig_box } ] }%
  \box_use_drop:N \l__pfig_box
  \tag_mc_end:
  \tag_struct_end:
}
\ExplSyntaxOff

% ----- links ------------------------------------------------------------
\usepackage{hyperref}
\hypersetup{colorlinks=true, urlcolor=AccentBlue, linkcolor=AccentBlue,
            pdfdisplaydoctitle=true}

% ----- poster vocabulary ------------------------------------------------
% Quiet section label: accent blue, tracked capitals. The label guides
% the eye; it does not decorate.
\NewDocumentCommand{\PosterSection}{m}{%
  \par\vspace{1.0em}%
  {\centering\SansSemibold\fontsize{26}{32}\selectfont\color{AccentBlue}%
   % LetterSpace pads the RIGHT of the last letter; a matching left kern
   % keeps the tracked caps optically centred.
   \hspace*{0.1em}\addfontfeatures{LetterSpace=10}\MakeUppercase{#1}\par}%
  \vspace{0.35em}}
% Run-in term label (sans emphasis; Inter has no italic).
\newcommand{\auterm}[1]{{\SansSemibold #1}}
\newcommand{\auhead}[1]{{\SansSemibold #1}}
% Side columns: calm. Centre column: slightly larger. One idea, one stage.
\newcommand{\SideBody}{\fontsize{24}{33.5}\selectfont}
\newcommand{\CenterBody}{\fontsize{25}{34.5}\selectfont}
\newcommand{\PosterCaption}{\fontsize{19}{26}\selectfont}
\newcommand{\PosterSmall}{\fontsize{20.5}{28}\selectfont}

% ============================  EDIT  ===================================
\newcommand{\DisplayTitle}{Calabi--Yau Manifolds}
\newcommand{\TechTitle}{What they are, why they exist, and why they matter}
\newcommand{\PosterAuthors}{Franklin Wang}
\newcommand{\PosterDept}{Department of Mathematics and Statistics, Auburn University}
\newcommand{\PosterURL}{https://github.com/Franklinwang72/auburn-university-latex-template}
\newcommand{\PosterURLShort}{github.com/Franklinwang72}
\newcommand{\PosterQRURL}{https://www.linkedin.com/in/franklin-wang-a14b7b327}
\newcommand{\PosterContact}{zzw0109@auburn.edu}
% PDF metadata: use REAL Unicode dashes here -- TeX ligatures like "--"
% are NOT applied inside PDF text strings and would show literally.
\hypersetup{pdftitle={Calabi–Yau Manifolds — What They Are, Why They Exist, and Why They Matter},
            pdfauthor={Franklin Wang}}
% NOTE: \DisplayTitle must fit ONE line at 100pt (about 40 characters).
% A wrapped title silently pushes the footer to a second page.
% ==========================  end EDIT  =================================
% Loud guard: warn if the poster no longer fits on one page.
\AtEndDocument{\clearpage\ifnum\value{page}>2
  \PackageWarningNoLine{poster-one}{Poster no longer fits on one page:
    shorten the title or body text, or reduce BodyH}\fi}

\begin{document}

% =========================  HEADER  ====================================
% Pure typography, centred: title, subtitle, author. The institution
% lives in the footer, next to the references, where it belongs.
\begin{center}
  {\SerifMedium\fontsize{100}{112}\selectfont \DisplayTitle\par}
  \vspace{0.45em}
  {\Serif\itshape\fontsize{42}{52}\selectfont\color{InkSoft} \TechTitle\par}
  \vspace{0.85em}
  {\fontsize{29}{38}\selectfont \PosterAuthors\;
   {\color{InkSoft}\fontsize{24}{38}\selectfont\textperiodcentered\; \PosterDept}\par}
\end{center}

\vspace{0.7em}
{\color{DividerWarm}\rule{\textwidth}{1.4pt}}

\vspace{1.0em}

% =========================  BODY  ======================================
% Three asymmetric columns. Fixed-height, stretchable minipages: \vfill
% distributes the whitespace so all three end exactly at the footer.
\newlength{\BodyH}\setlength{\BodyH}{50.4cm}
\begin{minipage}[t][\BodyH][s]{0.285\textwidth}
\SideBody

% ---------- LEFT: the prerequisites, built up for an undergraduate ----------
\PosterSection{What Is a Manifold?}
A \auterm{manifold} is a mathematician's word for a smooth shape --- of
any dimension. Its one defining trick: up close it looks \auterm{flat},
the way the ground looks flat from where you stand even though the
Earth is round. A circle, the skin of a doughnut, the fabric of
spacetime --- all are manifolds. The number of coordinates you need to
pin down a point is the \auterm{dimension}, and the shapes here are
\auterm{compact}: finite and edgeless, closed up like a sphere rather
than sprawling like an open plain.
\vfill
\PosterSection{Complex Dimensions}
Calabi--Yau manifolds run on \auterm{complex} coordinates. A complex
number $x + iy$ is a single arrow in the plane --- length and
direction bundled into one --- so one complex dimension already
carries two real ones. An $n$-complex-dimensional shape is therefore
$2n$-real-dimensional: a Calabi--Yau \auterm{threefold} $(n = 3)$ is a
\auterm{six}-dimensional world. Complex coordinates also unlock the
rigid, powerful toolbox of complex analysis.
\vfill
\PosterSection{Curvature: Do Triangles Behave?}
How can you tell a surface is curved without stepping outside it? Draw
a triangle. On a \auterm{sphere} its angles sum to \auterm{more} than
$180^\circ$; on a \auterm{saddle}, to \auterm{less}; only on a
\auterm{flat} sheet do they hit $180^\circ$ exactly. An ant that never
leaves the surface could discover its world is round this way ---
curvature is \auterm{intrinsic}, and it can change from point to point.
\begin{center}

\includegraphics[width=0.94\textwidth,
  alt={Three small surfaces in a row: a blue sphere with a triangle of
    three outward-bulging arcs, a flat gray sheet with a
    straight-edged triangle, and a tan saddle surface with a triangle
    whose sides pinch inward.}]{figures/curvature-trio.png}
\par\vspace{0.15em}
{\PosterCaption\color{InkSoft}
Sphere: more than $180^\circ$ \textperiodcentered\;
flat sheet: exactly $180^\circ$ \textperiodcentered\;
saddle: less than $180^\circ$.\par}
\end{center}
\vfill
\PosterSection{Ricci Curvature}
Past two dimensions, curvature splits into many independent pieces.
The \auterm{Ricci curvature} is the piece that controls \auterm{volume}.
Sprinkle a tiny ball of dust and let each speck drift along a straight
line: on a sphere the ball shrinks, on a saddle it swells, and when
$\operatorname{Ric} = 0$ it holds its volume exactly. Such a space is
genuinely curved yet volume-balanced --- which is precisely Einstein's
equation for a spacetime with no matter. \auterm{Ricci-flat spaces are
gravity's vacuums.}

\end{minipage}\hfill
%
\begin{minipage}[t][\BodyH][s]{0.37\textwidth}
\CenterBody

% ---------- CENTER: the definition, the theorem, the picture ----------
\PosterSection{What Is a Calabi--Yau Manifold?}
A \auterm{Calabi--Yau manifold} is a compact complex manifold with two
features. \auterm{(1) K\"ahler:} it carries both a way to measure
distance --- a metric --- and complex coordinates, fitting together as
smoothly as possible. \auterm{(2) A holomorphic volume form} $\Omega$:
a way to gauge complex volume, built from the coordinates, that
vanishes \emph{nowhere} (in bundle language, a trivial canonical
bundle $K_X \cong \mathcal{O}_X$). Feature (2) forces a purely
\auterm{topological} quantity to vanish --- the \auterm{first Chern
class} $c_1(X)$, a single whole-number invariant of the shape you can
compute without measuring anything. Calabi asked whether that cheap
test forces a Ricci-flat metric to exist.
\par\vspace{0.35em}
% The one lifted panel: soft blue-gray ground, hairline border,
% the statement set in Garamond -- a highlighted page in a book.
{\setlength{\fboxrule}{1.5pt}\setlength{\fboxsep}{0pt}%
\fcolorbox{ThmBorder}{ThmBG}{\begin{minipage}{\dimexpr\textwidth-3pt}
  \vspace{0.75em}\begin{center}\begin{minipage}{\dimexpr\textwidth-3.6em}
  % one quiet italic attribution -- no second caps voice next to the
  % MAIN RESULT label above the panel
  {\Serif\itshape\fontsize{26}{33}\selectfont\color{AccentBlue}%
   Theorem\ (Yau 1978 --- the Calabi conjecture)\par}
  \vspace{0.45em}
  {\Serif\fontsize{30}{42}\selectfont
  Let $(X,\omega)$ be a compact K\"ahler manifold with
  $c_1(X) = 0$ in $H^2(X,\mathbb{R})$. Then $X$ carries a
  \textit{unique} Ricci-flat K\"ahler metric $\tilde\omega$ with
  $[\tilde\omega] = [\omega]$:
  % the punchline at display scale (size set OUTSIDE the equation)
  {\fontsize{40}{47}\selectfont
  \begin{equation*}
    \operatorname{Ric}(\tilde\omega) \;=\; 0 .
  \end{equation*}}
  Every Calabi--Yau manifold admits Ricci-flat metrics ---
  exactly one in each K\"ahler class.\par}
  \end{minipage}\end{center}\vspace{0.75em}
\end{minipage}}}
\vfill
\auterm{How is it proved?}\ Fix an overall scale; every candidate
metric is then $\omega + i\,\partial\bar\partial\varphi$ for a single
function $\varphi$, and Ricci-flatness becomes one nonlinear PDE ---
the \auterm{complex Monge--Amp\`ere equation}
% size change OUTSIDE the environment: \fontsize inside math is a no-op
{\fontsize{29}{34}\selectfont
\begin{equation*}
  \bigl(\omega + i\,\partial\bar\partial\varphi\bigr)^{\,n}
  \;=\; e^{F}\,\omega^{n}
\end{equation*}}
--- whose unknown sits inside a determinant of its own second
derivatives. Yau's a~priori estimates force a solution to exist.
\vfill
\begin{center}
\PosterFigure{A three-quarter view of a complex-curve slice of the
  Fermat quintic Calabi-Yau threefold: interlocking petal-like sheets
  that flow from warm sand at the top through mist blue to deep blue at
  the base, with a fivefold pinwheel symmetry and a glossy finish.}{%
  \includegraphics[width=0.49\textwidth]{figures/calabi-yau-quintic.png}}
\par\vspace{0.25em}
% caption constrained to the figure width so the two read as one block
\begin{minipage}{0.63\textwidth}
{\PosterCaption\color{InkSoft}
A Calabi--Yau threefold has six real dimensions --- too many to draw,
but not too many to slice. Shown: the surface $z_1^5+z_2^5=1$, a
complex curve inside the Fermat quintic
$\,z_0^5+\dots+z_4^5=0\,$ in $\mathbb{P}^4$, projected to $3$-space;
its $25$ patches carry the standard Hanson parametrization.\par}
\end{minipage}
\end{center}

\end{minipage}\hfill
%
\begin{minipage}[t][\BodyH][s]{0.285\textwidth}
\SideBody

% ---------- RIGHT: examples a lay reader can hold, then the payoff ----------
\PosterSection{The Simplest One: A Donut}
In one complex dimension the definition selects exactly one shape: the
torus $\mathbb{C}/\Lambda$, the surface of a donut. Its volume form is
$dz$, inherited from the plane, and its Ricci-flat metric is the flat
geometry a sheet of paper keeps when rolled up in two directions ---
everything Yau's theorem promises, visible by hand.
\begin{center}
\includegraphics[width=0.36\textwidth,
  alt={A glossy blue-gray torus, the surface of a donut, with one
    vertical circle around the tube and one horizontal circle around
    the hole drawn in deep blue: the two directions in which the flat
    sheet was rolled up.}]{figures/cy-torus.png}
\par\vspace{0.1em}
{\PosterCaption\color{InkSoft}
The torus: a flat sheet rolled up in two directions.\par}
\end{center}
\vfill
\PosterSection{The Family Album}
One dimension up sit the celebrated \auterm{K3 surfaces}; in dimension
three, the \auterm{quintic}. A shape cut out by one degree-$d$
equation in complex projective space $\mathbb{P}^{n+1}$ is
Calabi--Yau exactly when $d = n+2$:
\begin{center}
\PosterSmall
% tag the first row as TH cells (screen readers announce column headers)
\tagpdfsetup{table/header-rows={1}}%
\begin{tabular}{c l l}
  \toprule
  \auhead{dim} & \auhead{Calabi--Yau} & \auhead{why $K_X \cong \mathcal{O}_X$} \\
  \midrule
  $1$ & elliptic curve $\mathbb{C}/\Lambda$ & $dz$ descends \\
  $2$ & K3 surface, e.g.\ $X_4 \subset \mathbb{P}^3$ & $4 = 2+2$ \\
  $3$ & quintic $X_5 \subset \mathbb{P}^4$ & $5 = 3+2$ \\
  \bottomrule
\end{tabular}
\end{center}
\vspace{0.2em}
The quintic alone moves in a $101$-parameter family of Ricci-flat
worlds.
\vfill
\PosterSection{Six Hidden Dimensions}
String theory needs spacetime to have \auterm{ten} dimensions --- four
we see, six curled up too small to notice. In 1985 Candelas, Horowitz,
Strominger and Witten showed that when the hidden six form a
Calabi--Yau threefold, the physics left in four dimensions keeps
supersymmetry, with Yau's Ricci-flat metric as its vacuum. If string
theory is right, a shape like the one at the centre of this poster
sits at every point of space.
\vfill
\PosterSection{A Shape and Its Mirror}
Calabi--Yau threefolds come in pairs: each $X$ has a \auterm{mirror}
$X^{\vee}$ with its two ``shape numbers'' swapped --- one counts ways
to resize the shape, the other ways to deform it. String theory cannot
tell the pair apart, and in 1991 Candelas--de~la~Ossa--Green--Parkes
used this to count rational curves on the quintic ($2875$ lines,
$609{,}250$ conics) by an integral on the mirror.
\vfill
\PosterSection{Takeaway}
% regular Garamond, ragged, unhyphenated: the theorem panel keeps its
% monopoly on typographic emphasis
{\Serif\fontsize{27}{38}\selectfont\raggedright
 \hyphenpenalty=10000\exhyphenpenalty=10000
A Calabi--Yau manifold passes one topological test, $c_1 = 0$ ---
and Yau's theorem rewards it with perfect geometry,
$\operatorname{Ric} = 0$: the meeting point of analysis, algebra, and
physics.\par}

\end{minipage}

\vspace*{\fill}

% =========================  FOOTER  ====================================
{\color{DividerWarm}\rule{\textwidth}{1.4pt}}
\vspace{0.6em}

% Footer on the body grid: the references block starts exactly at the
% centre column's left edge (0.315\textwidth); the QR block sits flush
% with the right margin.
\begin{minipage}[c]{0.315\textwidth}
  
\includegraphics[width=0.59\textwidth,
    alt={Auburn University formal logo}]{auburn/formal_horiz/onecolor_blue/auburn_formal_h_onecolor_blue_digital.png}
\end{minipage}%
\begin{minipage}[c]{0.545\textwidth}
  \fontsize{17.5}{24}\selectfont\color{InkSoft}
  \auhead{References:}\;
  E.~Calabi, Proc.\ ICM 1954 \textperiodcentered\;
  S.-T.~Yau, Comm.\ Pure Appl.\ Math.\ 31 (1978) \textperiodcentered\;
  P.~Candelas, X.~de~la~Ossa, P.~Green, L.~Parkes, Nucl.\ Phys.\ B 359 (1991) \textperiodcentered\;
  D.~Huybrechts, \textnormal{\Serif\itshape Complex Geometry}, Springer, 2005.\\
  \auhead{Contact:}\; \href{mailto:\PosterContact}{\PosterContact}\quad
  \auhead{Template:}\; \href{\PosterURL}{\PosterURLShort}\quad
  Set in EB Garamond, Inter \& STIX Two Math.
  \par
\end{minipage}\hfill
\begin{minipage}[c]{3.9cm}
  \centering
  % QR modules in main ink (a link would otherwise print accent blue,
  % costing scan contrast); keep the accent for the label only.
  \PosterFigure{QR code linking to the LinkedIn profile of \PosterAuthors\ at \PosterQRURL}%
               {\begingroup\hypersetup{urlcolor=InkMain}%
                \qrcode[height=3.9cm]{\PosterQRURL}\endgroup}\par
  \vspace{5pt}
  {\SansSemibold\fontsize{14}{19}\selectfont\color{AccentBlue}%
   \addfontfeatures{LetterSpace=8}CONNECT\par}
\end{minipage}

\end{document}
